%% abstract.tex
%% Abstract y palabras clave

%% abstract.tex
\begin{abstract}
El presente trabajo propone un modelo predictivo basado en \textit{Machine Learning} 
que funcione como \textbf{Sistema de Alertas Tempranas} para la detección de 
``Elefantes Blancos'' en contratos de infraestructura financiados por el Sistema 
General de Regalías (SGR). A partir de los registros de contratación disponibles 
en \textbf{SECOP II}, se integran variables contractuales y financieras-anticipos, 
modalidad de contratación, pluralidad de oferentes y ubicación— para predecir 
retrasos críticos e incumplimientos al momento de la adjudicación. Se entrenan y 
comparan clasificadores supervisados (\textit{Random Forest}, \textit{Regresión 
Logística}, \textit{XGBoost}), evaluados mediante validación cruzada estratificada 
y métricas robustas (AUC, F-score, precisión/recall). El sistema resultante busca 
transformar la auditoría reactiva en \textbf{vigilancia preventiva}, permitiendo a 
entidades de control y veedurías ciudadanas focalizar esfuerzos antes de que el 
detrimento patrimonial se materialice.
\end{abstract}

\begin{IEEEkeywords}
Elefantes Blancos, Sistema General de Regalías, Contratación pública, SECOP II, 
Aprendizaje automático, XGBoost, Random Forest, Alertas tempranas, 
Detección de anomalías, Auditoría preventiva
\end{IEEEkeywords}
